\section{Thermal calculation at inverse temperature \texorpdfstring{$49/4$}{49/4}}
\label{sec:thermal72}

We prove the $b=49/4$ table in Lemma~\ref{lem:thermal-inputs}. The computation
uses integer approximations to columns of a physical matrix exponential.
We specify the integer algorithm, bound its error, and give its exact
outputs. Throughout this section,
\[
 b=\frac{49}{4},\qquad J=280,\qquad Q=2^{55},\qquad 4\le n\le12.
\]
The twists are the identity and a rotation by \(\pi\) about the \(z\) axis.
The latter has the same partition function as \(P\), by the global rotation
described in Section~\ref{sec:model}.

\subsection{The physical matrix and its rational approximation}

In the weight basis, write a word as \(w=(w_0,\ldots,w_{n-1})\), with
\(w_i\in\{-1,0,1\}\). The total weight \(\sum_iw_i\) is conserved.
In each weight block let \(H\) be the restriction of \(H_n(g)\), and set
\begin{equation}\label{eq:thermal72-matrix}
 R=I-\frac{\frac32nI+H}{16},\qquad D=32R.
\end{equation}
By~\eqref{eq:seam-phase}, \(D\) is an integer matrix with diagonal
\[
 D_{w,w}=32-3n-2\sum_{i=0}^{n-1}w_{i-1}w_i.
\]
An allowed opposite unit step on the letters at \(i-1,i\) has entry
\(-2\), except that the entry is \(+2\) on the twisted seam \(i=0\).
These formulas specify every matrix entry; indices on words are modulo
\(n\). The unit step coefficient follows also directly from the two
ladder factors \(\sqrt2\) and the transverse prefactor \(1/2\).

Lemma~\ref{lem:energy} gives
\[
 0\le \frac32nI+H\le\frac52nI,\qquad
 \operatorname{spec}(R)\subseteq[1-5n/32,1]\subseteq[-7/8,1].
\]
Thus \(R\) is a self-adjoint contraction. Put
\[
 F=\exp\!\left[-\frac b2\left(\frac32nI+H\right)\right]
   =\exp[98(R-I)].
\]
We use the same symbols for the direct sums over all weight blocks.
On the full physical space,
\begin{equation}\label{eq:thermal72-trace}
 Z_n(b,g)=e^{x_n}\norm{F}_F^2,\qquad
 x_n=bn\left(\frac32-a\right),
\end{equation}
where \(\norm{\cdot}_F\) denotes the Frobenius norm and the squared norms
of the weight blocks are summed.

Define the rational probabilities
\begin{equation}\label{eq:thermal72-coefficients}
 p_j=\frac{v_j}{\sum_{\ell=0}^Jv_\ell},
 \qquad v_j=98^j\frac{J!}{j!},
 \qquad 0\le j\le J.
\end{equation}
They are Poisson weights of mean \(98\), conditioned to indices at most
\(J\). Since \(\norm{R}\le1\), discarding the Poisson tail and renormalizing
changes the exponential by at most twice the tail probability, by the
same coefficient estimate as in Fox--Glynn~\cite[Proposition~1]{FoxGlynn1988}:
\[
 \left\|F-\sum_{j=0}^Jp_jR^j\right\|
 \le 2e^{-98}\sum_{j=281}^{\infty}\frac{98^j}{j!}
 \le2^{-280}e^{98}<2^{-84}<Q^{-1}.
\]
For the second inequality, use \(2^j\ge2^{281}\) in the tail and
\(e^{-98}\sum_{j\ge0}196^j/j!=e^{98}\); the next uses \(e<4\).
The estimate includes the change of normalization.

\subsection{Integer columns and their error}

For a selected unit basis column \(e_r\), start with \(X=0\) and perform,
in the order \(j=J,J-1,\ldots,0\),
\begin{equation}\label{eq:thermal72-horner}
 X\ \longleftarrow\
 \left\lfloor\frac{DX}{32}\right\rfloor+c_je_r,
 \qquad c_j=\lfloor Qp_j\rfloor.
\end{equation}
The vector floor is coordinatewise. All quantities in this recursion are
integers. The coefficients \(c_j\) vanish for \(j>191\), so those initial
zero steps may be omitted without changing \(X\).

To bound rounding, compare with the ideal suffix
\[
 Y_j=Q\sum_{\ell=j}^Jp_\ell R^{\ell-j}e_r,\qquad \norm{Y_j}_2\le Q.
\]
In a block of dimension \(d_0\), one vector floor introduces Euclidean
error less than \(\sqrt{d_0}\), and one coefficient floor error less than
one. The contraction property therefore bounds the accumulated error
by \(281(\sqrt{d_0}+1)\). Set
\begin{equation}\label{eq:thermal72-errors}
 d_n=\lfloor\sqrt{3^n}\rfloor+1,\qquad
 E_n=281(d_n+1)+1,\qquad \mathcal E_n=d_nE_n.
\end{equation}
Including the preceding exponential-approximation error, every computed
column satisfies
\begin{equation}\label{eq:thermal72-column-error}
 \norm{X-QFe_r}_2\le E_n.
\end{equation}
The bound charges all \(281\) ideal steps, including the omitted steps
whose rational coefficients are nonzero but whose integer floors vanish.

The recursion can be evaluated exactly with signed \(64\)-bit integers.
Indeed every intermediate coordinate has absolute value at most
\[
 Q+281(d_n+1)<2Q,\qquad d_n\le730.
\]
The absolute sum of any row of \(D\) is at most
\[
 |32-3n|+2n+4n\le76.
\]
The first two terms bound the diagonal, and at most two hops per bond
give the last term. Consequently every product and every partial sum in
a sparse matrix multiplication has absolute value less than
\[
 76\cdot2Q=5476377146882523136<2^{63}-1.
\]
Division by \(32\) and addition of \(c_j\le Q\) also fit; even their
conservative absolute bound is
\[
 \frac{76\cdot2Q}{32}+Q=207165582859042816<2^{63}-1.
\]
Starting with \(X=0\), these bounds first establish exactness of the next
integer operation; the contraction estimate then establishes the next
coordinate bound. Thus the induction does not assume the absence of
overflow that it is proving. All final squares and sums below use
arbitrary-precision integers.

\subsection{Representative columns}

Only one column from each of the following word orbits is needed.
For positive total weight, use the orbit under cyclic shifts and
reversal, and give its representative twice the orbit size. This
includes the uncomputed negative-weight block. In the zero-weight
block, include digit negation in the orbit itself and use its size
without an extra factor two. Choose the lexicographically least word
as representative. Set cardinalities handle periodic words and all
other nontrivial stabilizers. The multiplicities \(\nu_r\) satisfy
\[
 \sum_r\nu_r=3^n.
\]

Here are the exact symmetries underlying this reduction. Reversal
\(w\mapsto(w_{n-1},\ldots,w_0)\) fixes the seam as an unordered bond, and global
negation preserves every Hamiltonian entry. For the twisted chain,
a bare shift moves the seam. Left rotation by \(q\) moves it to
\((k-1,k)\), where \(k=-q\bmod n\). The diagonal unitary
\[
 G_k|w\rangle=(-1)^{\sum_{j=0}^{k-1}w_j}|w\rangle,\qquad G_0=I,
\]
flips the hopping signs exactly at the two boundary bonds of this arc.
It restores the original seam. The resulting signed shift commutes with
\(H\), so exact exponential columns at rotated words have equal norms.
For the untwisted chain no sign correction is needed.

Let \(X_r\) be the output of~\eqref{eq:thermal72-horner} at representative
\(r\), and define the integer
\begin{equation}\label{eq:thermal72-total}
 W_{n,g}=\sum_r\nu_r\sum_w X_r(w)^2.
\end{equation}
Rounding need not preserve the signed symmetries. To compare norms,
form abstract direct sums repeating each exact representative column
\(\nu_r\) times, and do the same for each computed column.
The first direct sum has squared norm \(Q^2\norm{F}_F^2\), and the second
has squared norm \(W_{n,g}\). By~\eqref{eq:thermal72-column-error},
\begin{equation}\label{eq:thermal72-frobenius}
 \left|\sqrt{W_{n,g}}-Q\norm{F}_F\right|
 \le\sqrt{3^n}\,E_n<\mathcal E_n.
\end{equation}
This proves the full trace estimate using representatives, without
requiring their rounded columns to transform equivariantly.

\subsection{Exact totals and rational enclosures}

The integer algorithm~\eqref{eq:thermal72-coefficients},
\eqref{eq:thermal72-horner}, and~\eqref{eq:thermal72-total} gives the
following totals. The second column uses no twist; the third uses the
\(\pi\) twist.
\begin{center}
\begingroup
\footnotesize
\setlength{\tabcolsep}{4pt}
\begin{tabular}{r|rr}
\(n\)&\(W_{n,I}\)&\(W_{n,P}\)\\ \hline
4&1298092848946206529211955779616040&683496640248050713756440676071\\
5&9095379218684502088743241789&9205948219074287936685289758037\\
6&11971284263768914176431179720622&266573596119930013833919176645\\
7&15044150708248976451784593468&534330762492192297990602166370\\
8&386135473475360677201628357188&42459883501386986922766507327\\
9&4866417737994347551816846903&37716485750353424806876104310\\
10&19772960374215472827775892486&5030484777222468978572235564\\
11&841347687722481981930750335&2928767918985354021900689432\\
12&1266666171463993322695232916&524785463625387865619045258
\end{tabular}
\endgroup
\end{center}

We finish with a finite rational enclosure, specifying every bound used
to turn these integers into the thermal table at $b=49/4$. Let
\[
 r_{n,g}=\lfloor\sqrt{W_{n,g}}\rfloor,\qquad
 s_n=\sum_{j=0}^{100}\frac{x_n^j}{j!},\qquad
 t_n=\frac{x_n^{100}}{100!}.
\]
Since \(0\le x_n<15\), successive omitted terms have ratio at most
\(15/101\), and the tail is at most \(15t_n/86<t_n\). Hence
\(s_n\le e^{x_n}\le s_n+t_n\).
Equations~\eqref{eq:thermal72-trace} and~\eqref{eq:thermal72-frobenius}
give
\begin{equation}\label{eq:thermal72-enclosure}
 \begin{gathered}
 L_{n,g}=\frac{s_n\max(0,r_{n,g}-\mathcal E_n)^2}{Q^2},\qquad
 U_{n,g}=\frac{(s_n+t_n)(r_{n,g}+1+\mathcal E_n)^2}{Q^2},\\
 L_{n,g}\le Z_n(b,g)\le U_{n,g}.
 \end{gathered}
\end{equation}
All quantities in these endpoints are rational or integer, with the
integer square root determined by \(r^2\le W<(r+1)^2\).
Writing \(C_{n,g}\) for the integer centers in the $b=49/4$ table of
Lemma~\ref{lem:thermal-inputs}, direct integer and rational evaluation yields
\[
 L_{n,g}>\frac{C_{n,g}-30}{10^6},
 \qquad
 U_{n,g}<\frac{C_{n,g}+30}{10^6}
 \qquad(4\le n\le12,\ g=I,P).
\]
Every one of these strict inequalities has margin greater than
\(13/10^6\). The displayed algorithm, totals, and enclosures give a
reproducible finite verification of all eighteen enclosures, and
complete the proof of the thermal table at $b=49/4$.
